\documentclass[12pt]{article}
\usepackage{blog}

% ============================================================
% Blog Metadata
% ============================================================
\blogtitle{Understanding the Fast Fourier Transform (快速傅里叶变换)}
\blogdate{2025-12-01}
\blogtags{math, algorithms, signal-processing}
\bloglang{en}
\blogsummary{An introduction to the Fast Fourier Transform (FFT) --- from the DFT definition to the Cooley--Tukey algorithm, with Python code and complexity analysis.}

\begin{document}
\makeblogtitle

\section{Introduction}

computational mathematics and signal processing. It converts a sequence of $N$ complex
numbers into another sequence of $N$ complex numbers, revealing the frequency content of
the original signal.

The \textbf{Discrete Fourier Transform} (DFT)\footnote{See \textit{The Scientist and Engineer's Guide to Digital Signal Processing} by Steven W. Smith.} is one of the most important tools in computational mathematics and signal processing. It converts a sequence of $N$ complex numbers into another sequence of $N$ complex numbers, revealing the frequency content of the original signal.

\textcolor{blue}{DFT is the foundation of modern signal analysis.}

For a historical perspective, see \cite{cooley1965fft}.

然而，直接计算 DFT 需要 $O(N^2)$ 次运算，对于大规模数据来说效率太低。1965年，
Cooley 和 Tukey 发表了\textbf{快速傅里叶变换}（FFT）算法，将复杂度降低到
$O(N \log N)$，使得频谱分析在实际工程中变得可行。

This article walks through the mathematical foundation of the DFT, derives the radix-2
Cooley--Tukey FFT algorithm, and provides a Python implementation.

\section{The Discrete Fourier Transform}\label{sec:dft}

\begin{definition}[Discrete Fourier Transform]
Given a sequence $x_0, x_1, \ldots, x_{N-1}$ of complex numbers, its DFT is the sequence
$X_0, X_1, \ldots, X_{N-1}$ defined by:
\[
  X_k = \sum_{n=0}^{N-1} x_n \cdot e^{-i \, 2\pi k n / N}, \quad k = 0, 1, \ldots, N-1
\]
\end{definition}

We often write $\omega_N = e^{-2\pi i / N}$, the \emph{primitive $N$-th root of unity},
so the transform becomes:
\[
  X_k = \sum_{n=0}^{N-1} x_n \, \omega_N^{kn}
\]

The inverse DFT recovers the original sequence:
\[
  x_n = \frac{1}{N} \sum_{k=0}^{N-1} X_k \, \omega_N^{-kn}
\]

\subsection{Key Properties}

The DFT satisfies several important properties:

\begin{itemize}
  \item \textbf{Linearity}: $\mathrm{DFT}(\alpha x + \beta y) = \alpha\,\mathrm{DFT}(x) + \beta\,\mathrm{DFT}(y)$
  \item \textbf{Parseval's theorem}: $\sum_{n} |x_n|^2 = \frac{1}{N} \sum_{k} |X_k|^2$
  \item \textbf{Convolution theorem}: pointwise multiplication in the frequency domain corresponds
        to circular convolution in the time domain:
        \[
          \mathrm{DFT}(x * y) = \mathrm{DFT}(x) \cdot \mathrm{DFT}(y)
        \]
  \item \textbf{Shift property}: 时域中的移位对应频域中的相位旋转。若 $y_n = x_{n-m}$，则
        $Y_k = \omega_N^{mk} X_k$。
\end{itemize}

\section{The Cooley--Tukey FFT Algorithm}\label{sec:fft}

The key insight of the FFT is to exploit the symmetry and periodicity of $\omega_N$.

\begin{theorem}[Danielson--Lanczos Lemma]
An $N$-point DFT (where $N$ is even) can be decomposed into two $N/2$-point DFTs:
\[
  X_k = \underbrace{\sum_{m=0}^{N/2-1} x_{2m} \, \omega_{N/2}^{km}}_{E_k}
      + \omega_N^k \underbrace{\sum_{m=0}^{N/2-1} x_{2m+1} \, \omega_{N/2}^{km}}_{O_k}
\]
where $E_k$ is the DFT of the even-indexed elements and $O_k$ is the DFT of the
odd-indexed elements.
\end{theorem}

This gives us the \textbf{butterfly operation}:
\begin{align}
  X_k       &= E_k + \omega_N^k \, O_k \\
  X_{k+N/2} &= E_k - \omega_N^k \, O_k
\end{align}

Since $E_k$ and $O_k$ are periodic with period $N/2$, we only need to compute them for
$k = 0, 1, \ldots, N/2 - 1$. 通过递归地应用这一分解，我们可以将 $N$ 点 DFT 的计算
分解为 $\log_2 N$ 层蝶形运算，每层包含 $N/2$ 次蝶形操作。

\subsection{Complexity Analysis}

\begin{table}[H]
\centering
\caption{Comparison of DFT and FFT computational complexity}
\begin{tabular}{lccc}
\toprule
\textbf{Algorithm} & \textbf{Multiplications} & \textbf{Additions} & \textbf{Total} \\
\midrule
Naive DFT          & $N^2$                     & $N(N-1)$           & $O(N^2)$ \\
Radix-2 FFT        & $\frac{N}{2}\log_2 N$    & $N \log_2 N$       & $O(N \log N)$ \\
\bottomrule
\end{tabular}
\end{table}

For a concrete example, consider $N = 2^{20} \approx 10^6$:

\begin{itemize}
  \item Naive DFT: $\sim 10^{12}$ operations
  \item FFT: $\sim 10^7$ operations
  \item Speedup factor: $\sim 10^5$
\end{itemize}

\section{Python Implementation}\label{sec:python}

Below is a recursive radix-2 FFT implementation in Python.

\begin{lstlisting}[language=Python, caption={Recursive radix-2 FFT in Python}]
import numpy as np

def fft(x):
    """Compute the FFT of sequence x (length must be a power of 2)."""
    N = len(x)
    if N == 1:
        return x

    # Split into even and odd
    even = fft(x[0::2])
    odd  = fft(x[1::2])

    # Twiddle factors
    T = np.exp(-2j * np.pi * np.arange(N // 2) / N)

    # Butterfly
    return np.concatenate([
        even + T * odd,
        even - T * odd
    ])

# Example usage
if __name__ == "__main__":
    # Generate a signal: 50 Hz + 120 Hz
    fs = 1024           # Sampling rate
    t = np.arange(fs) / fs
    signal = np.sin(2 * np.pi * 50 * t) + 0.5 * np.sin(2 * np.pi * 120 * t)

    # Compute FFT
    spectrum = fft(signal)
    freqs = np.arange(fs) * fs / fs
    magnitudes = np.abs(spectrum) / fs

    print(f"Peak frequencies: {freqs[np.argsort(magnitudes)[-4:]]} Hz")
\end{lstlisting}

We can verify our implementation against NumPy's built-in FFT:

\begin{lstlisting}[language=Python, caption={Verification against NumPy}]
x = np.random.random(1024)
assert np.allclose(fft(x), np.fft.fft(x))
print("FFT implementation verified!")
\end{lstlisting}

\section{Applications}\label{sec:apps}

The FFT has far-reaching applications across many domains:

\begin{enumerate}
  \item \textbf{Signal processing}: spectral analysis, filtering, compression (e.g., MP3, JPEG)
  \item \textbf{Polynomial multiplication}: multiplying two degree-$n$ polynomials in $O(n \log n)$
        instead of $O(n^2)$
  \item \textbf{Large integer multiplication}: the Schönhage--Strassen algorithm uses FFT to
        multiply $n$-digit integers in $O(n \log n \log \log n)$
  \item \textbf{Partial differential equations}: 谱方法利用 FFT 在频域中高效求解偏微分方程，
        在流体力学和量子力学模拟中广泛使用
  \item \textbf{Convolution}: fast computation of convolutions via the convolution theorem,
        used in deep learning (convolutional neural networks)
\end{enumerate}

\section{Conclusion}

For more on the DFT, refer back to Section~\ref{sec:dft}. For the FFT algorithm, see Section~\ref{sec:fft}. For code, see Section~\ref{sec:python}. For real-world uses, see Section~\ref{sec:apps}.

\begin{thebibliography}{9}
\bibitem{cooley1965fft}
J. W. Cooley and J. W. Tukey, "An algorithm for the machine calculation of complex Fourier series," \textit{Mathematics of Computation}, vol. 19, no. 90, pp. 297--301, 1965.
\end{thebibliography}

FFT 是计算数学中最优美、最实用的算法之一。它将 DFT 的计算复杂度从 $O(N^2)$ 降低到
$O(N \log N)$，使得大规模频谱分析成为可能。From signal processing to number theory,
from image compression to solving PDEs, the FFT remains an indispensable tool in the
modern computational toolkit.

The key idea --- \emph{divide and conquer via the symmetry of roots of unity} --- is
both mathematically elegant and practically powerful. Understanding the FFT provides
deep insight into the interplay between the time domain and the frequency domain, a
duality that lies at the heart of much of applied mathematics.

\end{document}
